%% generated by tools/make_tables.py -- do not edit
\begin{table}[t]
\centering
\caption{六项验证实验的汇总。``对照量'' 一栏给出该项所对照的闭式解或参考实现，``偏差'' 一栏为实测值与参照值之差（相对量以百分数或科学计数法给出）。判据一栏是本文采用的验收标准：几何量须在双精度舍入水平内与闭式一致，渲染量须在亚像素内一致，管线量须在量化台阶内不可区分。末列为相应的物理或数值说明；全部数值由 \protect\path{validate_pytrace.json} 与 \protect\path{validate_accum.json} 经 \protect\path{tools/make_tables.py} 生成。}
\label{tab:v6summary}
\scriptsize
\begin{tabular}{l p{0.18\linewidth} p{0.15\linewidth} p{0.09\linewidth} p{0.17\linewidth} p{0.20\linewidth}}
\toprule
\# & 对照量 & 关键实测 & 偏差 & 判据 & 备注 \\
\midrule
V1 & 临界冲击参数 $b_{c}=3\sqrt{3}\,M$ & $b_{\rm edge}=5.196152422707$ & $9.99\times 10^{-15}$ & $|\Delta b|/b_{c}<10^{-12}$ & 二分搜索在双精度舍入极限处闭合 \\
V2 & 近临界轨道守恒性（$b=5.196147418$，绕光子球 $15.05$\,rad） & $\max$ 残差 $2.62\times 10^{-6}$ & 较 $tol=10^{-2}$ 下降 $159\times$ & 残差在 $tol\le 3\times10^{-4}$ 处触底 & 触底值为机器精度，继续收紧容差无改善 \\
V3 & 阴影半径的像素映射（闭式 $f_{\rm px}\tan\psi_{\rm sh}$） & $R_{\rm px}$ 亚像素反解 & $6.48\times 10^{-12}$ & 亚像素 $<10^{-9}$；整数扫描 $<1$\,px & 扫描口径的 0.84\% 偏差是像素量化，非物理误差 \\
V4a & 屏幕纵坐标 $\to$ 冲击参数 $b$ 的闭式 & $b_{\rm tracer}=12.08060346$ & $4.44\times 10^{-16}$ & $<10^{-12}$ & 几何映射与着色器实现逐位一致 \\
V4b & FIDO 角半径 $\sin\psi=(3\sqrt{3}M/r_{0})$ $\sqrt{1-2M/r_{0}}$ & $r_{0}=4M$：$\psi=66.716^{\circ}$ & $+34.90\%$ & FIDO 口径为准；flat/chart 均为近似 & 同一行 chart 口径 $-12.05\%$；$r_{0}\ge26M$ 后两口径误差才降到 $5\%$ 内 \\
V5 & 逃逸半径处的出射方向重建 & $\max=0.0206$\,px @ $1000\times600$ & -- & $<0.05$\,px & 误差随 $b$ 单调增长；$R_{\rm esc}=140M$ 在此视场下足够 \\
V6 & 渐进累积收敛（$n=1\to1024$，逐位重复） & rms 降 $9.04$ 倍，$\propto n^{-0.29}$ & -- & 8-bit 亮度地板 $=0$；三块重复对逐位相同 & 窗口内高频残余：r2 降 $24.5\%$，post 降 $14.8\%$ \\
V6 & 幂等性对照（$\Delta\equiv0$） & 重复对最大差异 $0.0$ & $0$ & 同一配置重复渲染须逐位相同 & 冻结抖动时 rms $\equiv0$：这是幂等性校验，不是精度校验 \\
\bottomrule
\end{tabular}

\end{table}
